%!TEX root = ../template.tex
%%%%%%%%%%%%%%%%%%%%%%%%%%%%%%%%%%%%%%%%%%%%%%%%%%%%%%%%%%%%%%%%%%%%
%% chapter6.tex
%% NOVA thesis document file
%%
%% Chapter with lots of dummy text
%%%%%%%%%%%%%%%%%%%%%%%%%%%%%%%%%%%%%%%%%%%%%%%%%%%%%%%%%%%%%%%%%%%%

\typeout{NT FILE chapter6.tex} %

\chapter{Conclusão}
\label{conclusao}

Este projeto visa a implementação do \textit{Dittor}, um mecanismo de mitigação de \textit{Sybil attacks} na rede Tor, com recurso a um esquema de credenciais baseado no \textit{SASSI} e assinaturas \textit{SyRA}. A partir da descentralização das autoridades de emissão de credenciais do esquema, é possível criar DIDs que representem os utilizadores voluntários que correm nós Tor, garantindo a impossibilidade de várias identidades para uma só entidade de forma anónima.

Apesar dos vários mecanismos criados para mitigação de \textit{Sybil attacks}, não existe uma solução teórica com implementação testada que o concretize de forma objetiva. O objetivo deste projeto será trazer inovação nesta área, com recurso a tecnologias desenvolvidas no último ano para implementar uma solução que realize a mitigação destes ataques num paradigma de impossibilidade de ameaça, ao invés de prevenção dos mesmos.

O processo de implementação terá uma duração estimada de nove meses, desde o processo de investigação à obtenção de resultados, que deverão ser analisados de forma a produzir uma conclusão de uso útil para o \textit{Tor Project}. O apoio contínuo da equipa de \textit{Network Health} do Tor será fulcral no encaminhamento para os objetivos pretendidos, garantindo que a implementação será útil para uma futura mudança no mundo do Tor ou que levará a novas ideias para tornar esta rede mais segura.